\documentclass[11pt,a4paper]{article}
\usepackage[T1]{fontenc}
\usepackage{lmodern}
\usepackage[margin=23mm]{geometry}
\usepackage{amsmath,amssymb,amsthm,mathtools,booktabs,array}
\usepackage[dvipsnames]{xcolor}
\usepackage[colorlinks=true,linkcolor=MidnightBlue,urlcolor=MidnightBlue,citecolor=MidnightBlue]{hyperref}
\usepackage{microtype}
\usepackage{fancyhdr}
\pagestyle{fancy}
\fancyhf{}
\fancyhead[L]{\small Erd\H{o}s 708: partial results}
\fancyhead[R]{\small 23 November 2026}
\fancyfoot[C]{\thepage}
\setlength{\headheight}{14pt}
\setlength{\parskip}{4pt}
\setlength{\parindent}{0pt}
\newtheorem{theorem}{Theorem}[section]
\newtheorem{lemma}[theorem]{Lemma}
\newtheorem{corollary}[theorem]{Corollary}
\newtheorem{proposition}[theorem]{Proposition}
\theoremstyle{remark}
\newtheorem{remark}[theorem]{Remark}
\newcommand{\pos}[1]{\left(#1\right)_+}
\newcommand{\ind}{\mathbf 1}
\newcommand{\pp}{P^+}
\newcommand{\vp}{v_p}
\newcommand{\sm}{\operatorname{sm}}
\newcommand{\Om}{\Omega}
\newcommand{\HH}{\mathcal H}
\newcommand{\NN}{\mathbb N}
\newcommand{\ZZ}{\mathbb Z}
\newcommand{\RR}{\mathbb R}
\DeclareMathOperator{\lcm}{lcm}
\title{\textbf{Partial Results on Erd\H{o}s Problem 708}\\[4pt]
\large Divisor certificates, exact covers for rough source sets,\\
and a precise remaining obstruction}
\author{Research working note prepared for Denis Paliy\\
\small AI-assisted derivations with independent internal audits}
\date{23 November 2026}
\begin{document}
\maketitle
\begin{center}
\fcolorbox{MidnightBlue}{blue!3}{\parbox{0.92\linewidth}{
\textbf{Status.} The unrestricted inequalities \(g(n)\le 2n\) and
\(g(n)\le(2+o(1))n\) are \textbf{not proved here}. This note gives
unconditional theorems for explicit classes, including an exact \(2n\)
theorem when all source prime factors are at least \(M^\alpha\),
\(\alpha>1/(2e)\), for sufficiently large \(M\).
The proofs have been independently checked within this research session;
they have not been externally refereed or formalized. No priority claim
is made. Finite searches are not used as substitutes for proofs.
}}
\end{center}
\begin{abstract}
For a set \(A\) of \(n\) distinct integers greater than one and
\(M=\max A\), the product-cover problem asks for a small subset of every
interval of \(M\) consecutive integers whose product is divisible by
\(\prod A\). We construct a signed divisor certificate that proves Hall's
condition for the quotients \(a/P^+(a)\), under an explicit condition on
the prime support of those quotients. A prime-valuation repair then
produces at most \(2n\) distinct interval elements. The cofactor primes
are at most \(\sqrt M\), giving the rough-source exponent \(1/(2e)\).
A second, ordered-prefix certificate proves a weighted hinge inequality
and a stronger level-set comparison when reciprocal-prime mass is below
one with a controlled rounding term. We also record an asymptotic
version allowing fixed small primes, a fixed-support uniformity theorem,
and a complete finite CRT formulation for counterexample searches.
\end{abstract}
\begingroup
\small
\setcounter{tocdepth}{1}
\setlength{\parskip}{1pt}
\makeatletter
\renewcommand{\l@section}[2]{\@dottedtocline{1}{0em}{2em}{#1}{#2}}
\makeatother
\tableofcontents
\endgroup
\newpage

\section{Problem and separation of claims}

Let \(A\subset\ZZ_{\ge2}\) be a set of \(n\) distinct integers, put
\(M=\max A\), and write
\[
 I_x=\{x+1,\ldots,x+M\},\qquad x\in\ZZ_{\ge0}.
\]
Define
\[
 \tau(A,I_x)=\min\{|B|:B\subseteq I_x,\ \prod_{a\in A}a\mid\prod_{b\in B}b\},
 \qquad
 g(n)=\sup_{|A|=n,\ x\ge0}\tau(A,I_x).
\]
Each interval position may be selected only once. The stronger target is
\(g(n)\le2n\); the asymptotic target is \(g(n)\le(2+o(1))n\).
The historical lower bound has leading coefficient two. Primary modern
working sources and their limitations are listed in
Section~\ref{sec:refs}; the new derivations below are stated separately
from those inputs.

We shall repeatedly use the elementary counts
\[
 \left\lfloor\frac M d\right\rfloor
 \le N_I(d):=\#\{b\in I:d\mid b\}
 \le \left\lceil\frac M d\right\rceil
 \le \frac M d+1.
 \tag{1.1}
\]
Also
\[
 \prod A\mid M!\mid\prod I_x.
 \tag{1.2}
\]
The first divisibility uses distinctness of the sources. The second
follows from \(\prod I_x=M!\binom{x+M}{M}\).

\begin{center}
\small
\begin{tabular}{p{0.40\linewidth}p{0.50\linewidth}}
\toprule
Statement & Status in this note\\
\midrule
Unrestricted \(g(n)\le2n\) & Open; not established\\
Unrestricted \(g(n)\le(2+o(1))n\) & Open; not established\\
Explicit sparse cofactor-support criterion & Exact \(2n\), proved\\
All source primes \(\ge M^\alpha\), \(\alpha>1/(2e)\) & Exact \(2n\) for all sufficiently large \(M\), proved\\
Cofactor reciprocal mass \(\le1-\delta\), fixed \(\delta>0\) & \(2n+o(n)\) uniformly in this class, proved using a cited long-interval theorem\\
Ordered prime-support condition & Weighted hinge and level-set inequalities, proved\\
Any fixed finite source prime support \(P\) & \(2n+|P|-1\) for sufficiently large \(M\), proved\\
CRT optimization experiments & Search evidence only; no general upper bound\\
\bottomrule
\end{tabular}
\end{center}

\section{An exact \(2n\) theorem from cofactor matching}
\label{sec:exact}

For \(a\ge2\), let \(p(a)=P^+(a)\) be its largest prime factor and
\(d(a)=a/p(a)\). Let \(A_c\) be the composite members of \(A\). Define
\[
 P=\{p:p\mid a\text{ for some }a\in A_c\},\qquad
 Q=\{q:q\mid d(a)\text{ for some }a\in A_c\}.
 \tag{2.1}
\]
For a prime set \(R\), put
\[
 \pi_R(t)=\#\{p\in R:p\le t\},\qquad H_R(t)=\sum_{\substack{p\in R\\p\le t}}\frac1p.
 \tag{2.2}
\]
The empty-composite case is immediate: choose a multiple of each prime
source and take their union, using at most \(n\) positions.

\begin{theorem}[Finite cofactor-support criterion]
\label{thm:exact}
Suppose \(A_c\ne\varnothing\). If for every distinct quotient
\(d=d(a)\), \(a\in A_c\), with \(T=M/d\), one has
\[
 H_Q(T)+\frac{\pi_Q(T)+\pi_P(T)+1}{T}\le1,
 \tag{2.3}
\]
then every \(I_x\) contains a set \(B\) with \(|B|\le2n\) and
\(\prod A\mid\prod B\).
\end{theorem}

\begin{proof}
Give each composite source \(a\) all target neighbors \(b\in I_x\)
satisfying \(d(a)\mid b\). We verify Hall's condition.

Fix an arbitrary subset \(S\subseteq A_c\), and let
\(D=\{d(a):a\in S\}\) be its distinct quotient set. Every \(d\in D\)
is \(Q\)-smooth and
\[
 dP^+(d)\le M.
 \tag{2.4}
\]
Indeed, if \(d=a/p(a)\), then \(P^+(d)\le p(a)\) and \(dp(a)=a\le M\).
For \(b\in I_x\), define the signed certificate
\[
 K_D(b)=\sum_{d\in D}
 \left(\ind_{d\mid b}-
       \sum_{\substack{q\in Q\\q\le M/d}}\ind_{dq\mid b}\right).
 \tag{2.5}
\]
At most one summand can be positive. To see this, positivity for \(d\)
means \(d\mid b\) and no prime \(q\in Q\), \(q\le M/d\), divides \(b/d\).
Together with (2.4), this says
\[
 d=\sm_{Q\cap[2,M/d]}(b),
 \tag{2.6}
\]
where the smooth part includes the \emph{full} powers of the indicated
primes in \(b\). If \(d_1<d_2\), the prime set on the right of (2.6)
for \(d_1\) contains the one for \(d_2\). The corresponding smooth part
is therefore at least as large, contradicting \(d_1<d_2\) if both
were fixed points.

Within \(I_x\), outside \(N(S)=\bigcup_{d\in D}(I_x\cap d\ZZ)\),
all summands vanish.
Inside it, the possible positive summand equals one and all others
are nonpositive. Hence
\[
 K_D(b)\le\ind_{b\in N(S)}.
 \tag{2.7}
\]
Using (1.1) and \(T=M/d\),
\[
 \begin{split}
 |N(S)|
 &\ge\sum_{d\in D}\left[N_I(d)-
                \sum_{\substack{q\in Q\\q\le T}}N_I(dq)\right]\\
 &\ge\sum_{d\in D}
       \left[T(1-H_Q(T))-\pi_Q(T)-1\right]\\
 &\ge\sum_{d\in D}\pi_P(T).
 \end{split}
 \tag{2.8}
\]
For each quotient \(d\), every source in its fibre has the form \(dp\)
with a prime \(p\in P\), \(p\le M/d\). Distinct sources give distinct
such primes. Thus
\[
 |S|\le\sum_{d\in D}\pi_P(M/d).
 \tag{2.9}
\]
Equations (2.8)--(2.9) prove Hall's condition. Match all composite
sources injectively. Prime sources have quotient one, so they may be
assigned afterward to unused positions; \(n\le M-1\) guarantees room.
We obtain \(n\) distinct positions \(B_0=\{b_a:a\in A\}\) with
\[
 \prod_{a\in A}d(a)\mid\prod B_0.
 \tag{2.10}
\]

It remains to supply the removed prime units. If
\(m_p=\#\{a\in A:p(a)=p\}\), the current deficit in the \(p\)-valuation
relative to \(\prod A\) is at most \(m_p\). The sum of all positive
deficits is therefore at most \(\sum_p m_p=n\).
Whenever a deficit remains, (1.2) ensures that some \emph{unused}
interval position contains that prime: otherwise the full interval
would have the same deficient valuation. Add such a position.
Each addition decreases the total deficit by at least one.
At most \(n\) additions finish the cover, proving \(|B|\le2n\).
\end{proof}

\begin{corollary}[An exact integer test using the actual fibres]
\label{cor:finite}
Let \(r_A(d)=\#\{a\in A_c:d(a)=d\}\). It is enough to check, for every
distinct composite-source quotient \(d\),
\[
 \left\lfloor\frac M d\right\rfloor
 -
 \sum_{\substack{q\in Q\\q\le M/d}}
 \left\lceil\frac M{dq}\right\rceil
 \ \ge\ r_A(d).
 \tag{2.11}
\]
This is a sufficient certificate valid for \emph{every} shift, not a
test of a single interval.
\end{corollary}
\begin{proof}
Use the exact floor/ceiling version of (2.8) and
\(|S|\le\sum_{d\in D}r_A(d)\).
\end{proof}

\begin{corollary}[A simple support condition]
\label{cor:cardinality}
Let \(P_0\) consist of \(k\) primes, all at least \(3k+1\).
Every finite set of distinct \(P_0\)-smooth integers satisfies the
exact \(2n\) conclusion, for every interval.
\end{corollary}
\begin{proof}
In (2.3), use \(H_Q(T)\le k/y\),
\(\pi_Q(T)+\pi_P(T)+1\le2k+1\), and \(T\ge y=\min P_0\).
The left side is at most \((3k+1)/y\le1\).
\end{proof}

For example, the 21 primes from 101 through 199 satisfy this condition
since \(101>3\cdot21+1=64\). Arbitrarily high powers are allowed.
This is a theorem about all finite source sets on that support, not
only squarefree triple products.

\section{Exact covers for rough source sets}
\label{sec:rough}

\begin{lemma}[The cofactor primes are small]
For \(Q\) in (2.1), one has \(Q\subseteq[2,\sqrt M]\).
\end{lemma}
\begin{proof}
If \(q\mid a/p(a)\), then \(q\le p(a)\) and \(qp(a)\mid a\).
Consequently \(q^2\le qp(a)\le a\le M\).
This includes the case \(q=p(a)\), where the largest prime is repeated.
\end{proof}

\begin{theorem}[Rough-source exponent]
\label{thm:rough}
Fix \(\alpha>1/(2e)\). There exists \(M_\alpha\) such that, whenever
\(M\ge M_\alpha\) and every prime divisor of every \(a\in A\) is at
least \(M^\alpha\), every interval of length \(M\) has a product cover
using at most \(2|A|\) distinct elements.
\end{theorem}
\begin{proof}
The case \(\alpha>1/2\) has only prime sources. At \(\alpha=1/2\),
the possible composite boundary cases satisfy the criterion directly
for large \(M\). Thus assume \(1/(2e)<\alpha<1/2\).
The preceding lemma gives
\[
 Q\subseteq[M^\alpha,M^{1/2}].
 \tag{3.1}
\]
The prime reciprocal asymptotic implies
\[
 H_Q(\infty)\le
 \sum_{M^\alpha\le q\le M^{1/2}}\frac1q
 =\log\frac1{2\alpha}+o(1)<1.
 \tag{3.2}
\]
Furthermore the elementary estimate \(\pi(t)=O(t/\log t)\) gives
\[
 \sup_{T\ge M^\alpha}\frac{\pi_Q(T)+\pi_P(T)+1}{T}
 \le\sup_{T\ge M^\alpha}\frac{2\pi(T)+1}{T}=o(1).
 \tag{3.3}
\]
There is a fixed positive gap between \(\log(1/(2\alpha))\) and one.
For sufficiently large \(M\), (3.2)--(3.3) imply (2.3).
Apply Theorem~\ref{thm:exact}.
\end{proof}

The constant is
\[
 \frac1{2e}=0.1839397205857\ldots.
 \tag{3.4}
\]
Equation (3.2) requires the \(o(1)\) form of the prime reciprocal
asymptotic, not merely an \(O(1)\) estimate. It follows, for instance,
by subtracting the logarithmic Mertens product formula at
\(M^{1/2}\) and \(M^\alpha\), and using
\(\log(1-1/p)=-1/p+O(p^{-2})\); see \cite{TaoMertens}.
No explicit numerical value of \(M_\alpha\) is asserted here.

\section{An asymptotic class allowing small primes}
\label{sec:asymclass}

\begin{theorem}
\label{thm:asymclass}
Fix \(\delta>0\). For all source sets whose cofactor support \(Q\)
satisfies
\[
 \sum_{q\in Q}\frac1q\le1-\delta,
 \tag{4.1}
\]
one has uniformly
\[
 \tau(A,I_x)\le2n+o_\delta(n).
 \tag{4.2}
\]
The minimum prime of \(Q\) need not tend to infinity.
\end{theorem}
\begin{proof}
Choose a constant \(T_0=T_0(\delta)\ge2\) such that
\[
 \frac{2\pi(T)+1}{T}\le\delta\qquad(T\ge T_0).
 \tag{4.3}
\]
Call a composite source good if \(M/d(a)\ge T_0\), and bad otherwise.
For all good quotient vertices the proof of Theorem~\ref{thm:exact}
applies, using (4.1)--(4.3); match them injectively. Prime sources may
again be included by their full neighborhoods.

If a source is bad, then
\[
 P^+(a)\le M/d(a)<T_0.
 \tag{4.4}
\]
It is therefore composed entirely of the fixed finite collection of
primes below \(T_0\). The number \(b\) of distinct bad sources is at most
\[
 b\le\prod_{p<T_0}
       \left(1+\left\lfloor\frac{\log M}{\log p}\right\rfloor\right)
 =O_\delta((1+\log M)^{\pi(T_0)}).
 \tag{4.5}
\]
Match the good cofactors and then repair all valuation deficits.
The removed largest prime units of the good sources cost at most their
number. Supply the \emph{entire} prime factor demand of the bad sources
by the same unused-position repair rule. The latter cost is at most
\(\sum_{\text{bad }a}\Omega(a)\le b\log_2M\).
The full-interval divisibility (1.2) ensures every repair exists.
Thus
\[
 \tau(A,I_x)\le
 2n+O_\delta((1+\log M)^{\pi(T_0)+1}).
 \tag{4.6}
\]
This does not by itself give a bound in \(n\), since \(M\) is unbounded.
Use the established long-interval result
\(M\ge16n^2\Rightarrow\tau(A,I_x)\le2n\) from \cite{Wang}.
In the remaining range \(M<16n^2\), the error in (4.6) is a fixed
power of \(\log n\), hence \(o_\delta(n)\).
The older cubic long-interval theorem from \cite{Chen} also suffices
for this last reduction.
\end{proof}

The fixed margin \(\delta\) is essential to the stated uniform
conclusion. We have not proved (4.2) when the reciprocal mass is
unrestricted, or when its distance from one decays without control.

\subsection{An arbitrary fixed collection of small prime factors}

The exact matching argument can absorb a fixed set of small primes.
This is a source-set statement; the number of such primes need not
be two.

\begin{proposition}[Fixed small-prime part]
\label{prop:fixedanchors}
Let \(E\) and \(R\) be disjoint finite prime sets, with every member
of \(E\) smaller than \(y=\min R\). Define
\[
 \theta_E=\prod_{q\in E}(1-1/q),\quad C_E=2^{|E|},\quad
 \rho_y(R)=\sup_{T\ge y}\frac{\pi_R(T)+1}{T},\quad
 \sigma_y(R)=\sup_{T\ge y}\frac{\pi_R(T)}{T}.
 \]
If
\[
 \theta_E\left(1-\sum_{r\in R}\frac1r\right)
 \ge C_E\rho_y(R)+\sigma_y(R),
 \tag{4.7}
\]
then every distinct source set supported on \(E\cup R\), each of whose
members has a prime factor in \(R\), has an exact \(2n\) cover.
\end{proposition}
\begin{proof}
For a composite source the removed largest prime belongs to \(R\).
Thus every quotient \(d\) used by a Hall subset has \(T=M/d\ge y\).
Let \(L_E=\prod_{q\in E}q\) and replace (2.5) by
\[
 \sum_{\substack{d\in D\\d\mid b}}
 \ind_{\gcd(b/d,L_E)=1}
 \left(1-\sum_{\substack{r\in R\\r\le M/d}}\ind_{r\mid b/d}\right).
 \tag{4.8}
\]
A positive summand forces \(d\) to be the full smooth part of \(b\)
on \(E\cup(R\cap[2,M/d])\). These prime sets decrease as \(d\)
increases, so again at most one summand is positive.
The expression is therefore at most the neighborhood indicator.

Inclusion--exclusion over the fixed set \(E\) counts integers
coprime to \(L_E\) in a real interval of length \(T\) as
\(\theta_ET+O(C_E)\), with absolute error at most \(C_E\).
Each additional \(r\)-divisibility condition has main term
\(\theta_ET/r\) and the same error, since \(r\notin E\).
Thus the contribution of \(d\) to the neighborhood lower bound
is at least
\[
 T\theta_E\left(1-\sum_{r\in R}\frac1r\right)
 -C_E(\pi_R(T)+1)\ge\pi_R(T).
 \]
The source fibre over \(d\) has at most \(\pi_R(T)\) members.
Hall's theorem, separate quotient-one matching for prime sources,
and the same valuation repair finish the proof.
\end{proof}

For fixed \(E\) and fixed \(\delta>0\), condition (4.7) holds uniformly
when \(\sum_{r\in R}1/r\le1-\delta\) and \(y\) is sufficiently large.
If anchor-only sources are also allowed, their number is bounded
by a fixed power of \(\log M\). Supplying their full prime-unit demands,
as in Theorem~\ref{thm:asymclass}, gives the asymptotic \(2n+o(n)\)
conclusion in this enlarged class. Neither \(E\) nor its cardinality
may silently grow in that uniform assertion.

\section{Weighted inequalities from ordered prime prefixes}
\label{sec:hinge}

The preceding exact theorem concerns the full prime factorization of
the sources. Weighted hinge statements have a different scope: an
integer may contain prime factors outside the weighted support.
The two settings must not be conflated.

Let \(P\) be a nonempty finite prime set and \(y=\min P\); the empty
support is trivial. For \(p\in P\), let
\(u_p(e)\), \(e\ge0\), be nondecreasing, with \(u_p(0)=0\) and
\(0\le u_p(e)\le1\). Put
\[
 U(a)=\sum_{p\in P}u_p(v_p(a)),\qquad
 m_p=\sum_{j\ge1}\frac{u_p(j)-u_p(j-1)}{p^j}\le\frac1p,
 \tag{5.1}
\]
and
\[
 \rho_y(P)=\sup_{T\ge y}\frac{\pi_P(T)+1}{T}.
 \tag{5.2}
\]
For a finite ordered support \(p_1<\cdots<p_k\), this is exactly
\(\max_i(i+1)/p_i\).

\begin{theorem}[Ordered level comparison]
\label{thm:levels}
If
\[
 H_P(\infty)+\rho_y(P)\le1,
 \tag{5.3}
\]
then, for every real \(s\ge0\), every \(M\), and every interval \(I\)
of \(M\) consecutive positive integers,
\[
 \#\{a\le M:U(a)>s+1\}
 \le
 \#\{b\in I:U(b)>s\}.
 \tag{5.4}
\]
\end{theorem}

\begin{proof}
For a \(P\)-smooth \(d>1\), write \(p(d)=P^+(d)\) and
\(e(d)=v_{p(d)}(d)\). Let \(\mathcal F_s\) contain precisely those \(d\)
such that
\[
 dp(d)\le M,\qquad
 U(d)>s\ge U(d/p(d)^{e(d)}).
 \tag{5.5}
\]
For a source \(a\) counted on the left of (5.4), form the complete
prime-power prefixes of its \(P\)-smooth part, in increasing prime
order. Let \(d\) be the first prefix with weight above \(s\).
The last added coordinate has weight at most one, so \(U(d)\le s+1\).
The prefix is proper, since \(U(a)>s+1\); a strictly larger weighted
prime remains outside \(d\). In particular \(dp(d)\le M\), so
\(d\in\mathcal F_s\). Counting that outside prime gives
\[
 \#\{a\le M:U(a)>s+1\}
 \le
 M\sum_{d\in\mathcal F_s}\frac1d
       \sum_{\substack{q\in P\\q>p(d)}}\frac1q.
 \tag{5.6}
\]
This estimate follows by the union bound over outside primes and
\(\lfloor M/(dq)\rfloor\le M/(dq)\). Prime factors of \(a\) outside
\(P\) are harmless.

Use the ordered certificate
\[
 K_s(b)=\sum_{d\in\mathcal F_s}
 \left(\ind_{d\mid b}-
       \sum_{\substack{q\in P\\q\le p(d)}}\ind_{dq\mid b}\right).
 \tag{5.7}
\]
A positive summand must be an actual complete prime-power prefix of
the \(P\)-smooth part of \(b\). There is at most one such prefix that
crosses the threshold \(s\). Every other summand is nonpositive, so
\[
 K_s(b)\le\ind_{U(b)>s}.
 \tag{5.8}
\]
For \(T=M/d\ge p(d)\ge y\), interval counting gives
\[
 N_I(d)-\sum_{\substack{q\in P\\q\le p(d)}}N_I(dq)
 \ge
 T\left(1-H_P(p(d))-\rho_y(P)\right).
 \tag{5.9}
\]
By (5.3), the last bracket is at least
\(\sum_{q>p(d),\,q\in P}1/q\). Sum (5.8)--(5.9) and compare with
(5.6).
\end{proof}

\begin{corollary}[Bounded-profile hinge]
Under (5.3),
\[
 \sum_{a\le M}\pos{U(a)-2}
 \le \sum_{b\in I}\pos{U(b)-1}.
 \tag{5.10}
\]
\end{corollary}
\begin{proof}
Integrate (5.4) over \(s\in[1,\infty)\). All functions are bounded
for this fixed finite instance.
\end{proof}

\begin{proposition}[An ordered-increment hinge proof]
For (5.10), it is enough that for every \(p\in P\),
\[
 H_P(p)+\sum_{\substack{q\in P\\q>p}}m_q+\rho_y(P)\le1.
 \tag{5.11}
\]
\end{proposition}
\begin{proof}
For a list \(z_i\in[0,1]\) in prime order, let \(S_i=\sum_{j\le i}z_j\),
\(S=\sum z_i>2\), and
\[
 h_i=\pos{S_i-1}-\pos{S_{i-1}-1}.
 \]
Then \(0\le h_i\le1\), \(\sum h_i=S-1\), and
\[
 \sum_i h_i(S-S_i)=\sum_{i<j}h_ih_j\ge S-2.
 \tag{5.12}
\]
For the equality, every increment after a positive \(h_i\) equals
the corresponding \(z_j\). For the inequality, the multiaffine
function \(\sum t_i-\sum_{i<j}t_it_j\) is at most one on the unit
cube: at a vertex with \(r\) ones it is \(r-\binom r2\le1\).

For each \(P\)-smooth \(d\) with \(dp(d)\le M\), set
\[
 h_d=\pos{U(d)-1}-\pos{U(d/p(d)^{e(d)})-1}.
 \]
Apply (5.12) to the source prefixes. After summation, the left hinge
is at most
\[
 M\sum_d\frac{h_d}{d}\sum_{q>p(d)}m_q.
 \]
Use (5.7) with the coefficient \(h_d\) and without a crossing
restriction. The positive summands are genuine prefixes; their
nonnegative coefficients telescope to \(\pos{U(b)-1}\).
Equation (5.9), multiplied by \(h_d\), and (5.11) finish the proof.
\end{proof}

\subsection{Return to completely additive prime weights}

For any nonnegative weights \(w_p\), put
\[
 X_w(a)=\sum_{p\in P}w_pv_p(a),\quad
 u_p(e)=\min(w_pe,1),\quad t_p(e)=\pos{w_pe-1},
 \quad T(a)=\sum_p t_p(v_p(a)).
 \]
Then
\[
 \pos{X_w(a)-1}=T(a)+\pos{U(a)-1},\qquad
 \pos{X_w(a)-2}\le T(a)+\pos{U(a)-2}.
 \tag{5.13}
\]
Every \(t_p\) is nondecreasing from zero. Expanding its successive
increments into prime-power divisibility indicators and using (1.1)
shows \(\sum_{a\le M}T(a)\le\sum_{b\in I}T(b)\).
Thus (5.10) implies the raw hinge inequality
\[
 \sum_{a\le M}\pos{X_w(a)-2}
 \le\sum_{b\in I}\pos{X_w(b)-1}.
 \tag{\(H_2\)}
\]
In particular, for each fixed \(\delta>0\), \(H_2\) holds uniformly
for every finite support with sufficiently large minimum prime and
\(\sum_{p\in P}1/p\le1-\delta\).

\section{Adding small prime anchors}

\begin{proposition}[One arbitrary extra prime]
\label{prop:anchor}
Suppose \(P\) satisfies (5.3), and \(q\notin P\) is any prime.
Then \(H_2\) holds for all nonnegative raw weights supported on
\(P\cup\{q\}\).
\end{proposition}
\begin{proof}
First use bounded profiles. Add the \(q\)-profile to \(U\) and
telescope over its nonnegative increments
\([u_q(r-1),u_q(r)]\subseteq[0,1]\). At level \(r\), multiples of
\(q^r\) in the initial interval, after division, form
\(\{1,\ldots,N\}\), \(N=\lfloor M/q^r\rfloor\).
The target multiples, after division, form \(N\) or \(N+1\)
consecutive integers. For an integration value \(v\in[0,1]\), the
source hinge increment counts \(U>2-v\), whereas the target
increment counts \(U>1-v\). Theorem~\ref{thm:levels} with
\(s=1-v\ge0\), applied to any \(N\)-term target subinterval,
compares these increments. \(U(q^rk)=U(k)\) because \(q\notin P\).
Add the baseline (5.10), integrate, and then apply (5.13).
\end{proof}

The conclusion is a hinge inequality for the enlarged support,
not the level comparison (5.4). It therefore does not automatically
iterate to arbitrarily many extra primes.

\begin{theorem}[Two arbitrary smaller primes]
\label{thm:twoanchors}
Let \(E\) contain at most two primes, all less than \(y\), and let all
primes in \(R\) be at least \(y\). With \(\theta_E,C_E\) as above, the condition
\[
 \sum_{r\in R}\frac1r+\frac{C_E}{\theta_E}\rho_y(R)\le1
 \tag{6.1}
\]
implies the bounded-profile hinge inequality on \(E\cup R\), and
therefore raw \(H_2\) for all nonnegative weights on that support.
For \(|E|\le2\), the uniform constant \(C_E/\theta_E\le12\) suffices.
\end{theorem}
\begin{proof}
Use the ordered hinge increments \(h_d\) from the proof of (5.11),
retaining only positive coefficients and carriers with
\(T=M/d\ge\max(P^+(d),y)\). A positive increment requires at least
two distinct active prime coordinates. Consequently
\(p=P^+(d)\) is no smaller than every prime in \(E\).
Every nonzero residual source term in (5.12)
leaves a strictly larger prime from \(R\); it also leaves a cofactor
coprime to \(L_E\). The restriction \(T\ge y\) therefore loses no
source contribution.

Keep this coprimality condition in the source upper estimate.
For a fixed carrier,
\[
 \sum_{\substack{m\le T\\(m,L_E)=1}}U_{R,>p}(m)
 \le
 \theta_ET\sum_{\substack{r\in R\\r>p}}m_r
 +C_E\bigl(\pi_R(T)-\pi_R(p)\bigr).
 \tag{6.2}
\]
Indeed, expand each profile into nonnegative prime-power increments.
Only powers at most \(T\) contribute to the error, and the increments
of each prime sum to at most one.

For the target certificate use
\[
 K(b)=\sum_{\substack{d\text{ a carrier}\\d\mid b}} h_d
 \ind_{\gcd(b/d,L_E)=1}
 \left(1-\sum_{\substack{r\in R\\r\le p(d)}}\ind_{r\mid b/d}\right).
 \tag{6.3}
\]
Its positive terms are genuine complete prime prefixes; their
coefficients telescope, so \(K(b)\le\pos{U(b)-1}\).
The interval contribution for a carrier is at least
\[
 h_d\left[
 \theta_ET\left(1-\sum_{\substack{r\in R\\r\le p}}\frac1r\right)
 -C_E(\pi_R(p)+1)\right].
 \tag{6.4}
\]
Subtract \(h_d\) times the source upper estimate (6.2). The two error counts
combine to \(C_E(\pi_R(T)+1)\), while \(m_r\le1/r\).
The resulting difference is at least
\[
 h_d\theta_ET
 \left[1-\sum_{r\in R}\frac1r
             -\frac{C_E}{\theta_E}\rho_y(R)\right]\ge0.
 \]
Sum over \(d\) and apply (5.13).
\end{proof}

In particular \(H_2\) holds on \(\{2,3\}\cup R\), uniformly over all
weights and interval lengths, if \(\sum_{r\in R}1/r\le1-\delta\)
and \(\min R\) is sufficiently large. The proof does not automatically
absorb a third smaller prime: an early positive hinge increment
could leave another such prime outside, invalidating both the
coprimality condition and the cutoff \(T\ge y\).

\section{Uniform eventual \(H_2\) on a fixed prime support}

\begin{theorem}
\label{thm:fixed}
For every fixed finite prime set \(P\), there is \(M_0(P)<\infty\)
such that \(H_2\) holds for all \(M\ge M_0(P)\), all shifts, and all
weights \(w_p\in[0,1]\) supported on \(P\).
\end{theorem}
\begin{proof}
The zero-weight case is immediate. Write \(Q=\max P\) and
\(\rho=\max_{p\in P}w_p/\log p>0\). Let
\[
 C_P=\prod_{p\in P}\frac{1-1/p}{1-p^{-1/8}}.
 \]
Choose \(\rho_0>0\) so small that
\[
 Qe^{-1/(2\rho_0)}\le\tfrac12,\qquad
 \frac{32QC_P\rho_0}{7e}e^{-1/(4\rho_0)}\le1.
 \tag{7.1}
\]
For \(0<\rho\le\rho_0\), put \(t=7/(8\rho)\). The multiplicative
function \(e^{tX_w(n)}\) has a nonnegative divisor expansion, since
its prime-power successive increments are nonnegative. Thus
\[
 \sum_{a\le M}e^{tX_w(a)}
 \le M\prod_{p\in P}\frac{1-1/p}{1-e^{tw_p}/p}
 \le MC_P.
 \]
Since \(\pos{u-2}\le e^{t(u-2)}/(et)\),
\[
 L:=\sum_{a\le M}\pos{X_w(a)-2}
 \le M\frac{8\rho C_P}{7e}e^{-7/(4\rho)}.
 \tag{7.2}
\]
Also \(X_w(a)\le\rho\log a\); hence \(L=0\) if \(M\le e^{2/\rho}\).
Otherwise choose a prime attaining \(\rho\), and let
\(j=\lceil3/(2w_p)\rceil\), \(D=p^j\).
Then \(jw_p-1\ge1/2\) and
\[
 D\le Qe^{3/(2\rho)}\le M/2.
 \]
The interval contains at least \(M/(2D)\) multiples of \(D\), so
\[
 R:=\sum_{b\in I}\pos{X_w(b)-1}
 \ge\frac{M}{4Q}e^{-3/(2\rho)}.
 \tag{7.3}
\]
By (7.1), (7.2) is at most (7.3).

It remains to treat \(\rho\ge\rho_0\), uniformly in the weights.
Put \(k=|P|\), \(a_0=\rho_0\log2\),
\(T_0=\lceil2/a_0\rceil\), and \(\delta_0=Q^{-T_0}\).
Choose \(T\ge T_0\) such that
\[
 \sum_{p\in P}\frac1{p^T(p-1)}\le\delta_0/2,
 \]
and put \(Y(a)=\sum_pw_p\min(v_p(a),T)\).
This function has period \(R_0=\prod_pp^T\) and is bounded by \(kT\).
Some \(w_p\ge a_0\); on multiples of \(p^{T_0}\), \(Y\ge2\).
Therefore the period mean of
\(\pos{Y-1}-\pos{Y-2}\) is at least \(\delta_0\).
Either periodic hinge differs from its period mean times \(M\) by
at most \(R_0kT\). The prefix truncation error is at most
\(M\sum_p[p^T(p-1)]^{-1}\), while truncation only decreases the
target hinge. Consequently
\[
 L-R\le-\delta_0M/2+2R_0kT.
 \]
It is nonpositive for \(M\ge4R_0kT/\delta_0\).
\end{proof}

The constants depend on the actual support, not only its cardinality.
This theorem excludes amplification by keeping the prime set fixed;
it does not rule out counterexamples on growing supports.

\begin{corollary}[Integer covers on every fixed prime support]
\label{cor:fixedinteger}
Let \(P\) be a fixed nonempty finite prime set, \(k=|P|\). For
\(M\ge M_0(P)\), every set \(A\subseteq[2,M]\) of \(n\) distinct
\(P\)-smooth integers has, in every positive interval \(I\) of
length \(M\), a product cover with
\[
 \tau(A,I)\le2n+k-1.
 \tag{7.4}
\]
In particular, the upper bound \((2+o(1))n\) holds on each fixed
prime support, uniformly in the interval shift.
\end{corollary}
\begin{proof}
Consider the boxed fractional covering program (9.2), with its
\(k\) valuation constraints restricted to \(P\). Its dual weights
may be clipped to \([0,1]\), as explained after (9.3). By
Theorem~\ref{thm:fixed}, every dual objective is at most
\[
 \begin{split}
 \sum_{a\in A}X_w(a)-\sum_{b\in I}\pos{X_w(b)-1}
 &\le\sum_{a\in A}X_w(a)-\sum_{m\le M}\pos{X_w(m)-2}\\
 &\le\sum_{a\in A}\min\{X_w(a),2\}\le2n.
 \end{split}
\]
Thus the fractional optimum is at most \(2n\).
An optimal vertex has at most \(k\) coordinates strictly between
zero and one: otherwise a nonzero perturbation supported on those
coordinates could preserve all \(k\) valuation constraints, contrary
to extremality. Round these coordinates upward. The valuation
constraints remain satisfied. If there is a fractional coordinate,
the cost increase is strictly less than \(k\), so the integer cover
has fewer than \(2n+k\) positions, hence at most \(2n+k-1\).
If none is fractional, its cost is already at most \(2n\).
Finally, distinctness implies \(M\ge n+1\), so the threshold
\(M\ge M_0(P)\) is automatic for all sufficiently large \(n\).
\end{proof}

\section{A finite, complete CRT search model}

For fixed \(M\) and a prime \(p\le M\), let
\[
 E=\lfloor\log_pM\rfloor,\qquad d=p^E,\qquad Q_p=p^{E+1}>M.
 \]
For minimizing the target hinge there is no need for valuations above
\(E+1\): at most one interval position is divisible by \(Q_p\), and
one further \(p\)-adic digit can make its valuation exactly \(E+1\)
without changing any other position.

Write \(s=-x\bmod Q_p\). For each fixed \(s\bmod d\), the first \(E\)
valuation levels are fixed. Each next-digit choice with
\(s\in\{1,\ldots,M\}\) gives one high position; there are \(M\) such
choices in total. All empty choices within the same residue modulo
\(d\) give the same vector, so at most \(d\) additional representatives
are needed. Hence at most
\[
 \min(Q_p,M+d)\le2M
 \tag{8.1}
\]
valuation vectors per prime suffice. These are real interval
patterns: compatible choices for different primes are simultaneously
realized by the Chinese remainder theorem, with an extra digit used
to fix the possible high valuation.

Thus the search over all shifts and continuous weights is exactly a
finite mixed integer linear optimization problem for fixed \(M,P\).
The supplied code implements this model, coordinate optimization, and
rational re-evaluation of candidate witnesses. Floating-point solver
bounds are not exact proof certificates.

The completed local searches found no violation of \(H_2\).
They are not exhaustive in \(M\), support, phases, or weights.
Global mixed integer runs that timed out are expressly not described
as completed verifications. The more important outputs are the
proofs in Sections~\ref{sec:exact}--\ref{sec:hinge}.

Two concrete exact checks accompany the theorem:
\begin{itemize}
\item All \(1330=\binom{21}{3}\) products of three distinct primes
from 101 through 199 have \(M=7\,566\,179\). Both finite criteria
in Section~\ref{sec:exact} pass, and the minimum surplus in (2.11)
is 158. Hence every interval of this length has a cover using
at most 2660 positions. No graph with all \(M\) target vertices
was built for this certificate.
\item On the support \(\{11,13,17,19,23\}\), take all 35 cubic
monomials, the five squares, and the five primes. There are 45
distinct sources, \(M=12167\). In a length-\(M\) interval centered
at their least common multiple, the algorithm constructs a
45-position cofactor matching and makes 36 repairs, producing an
81-position cover within the 90-position bound. The complete
matching, repair trace, and valuation check are included in JSON.
\end{itemize}
The JSON can be re-verified with the Python standard library alone;
SciPy was used only to find the illustrative matching.

\section{What would finish the unrestricted asymptotic problem}
\label{sec:remaining}

For \(w_p\in[0,1]\), define
\[
 \Delta_c(M,x,w)=
 \sum_{a=1}^M\pos{X_w(a)-c}
 -
 \sum_{b\in I_x}\pos{X_w(b)-1}.
 \tag{9.1}
\]
The boxed fractional covering program is
\[
 L(A,I)=\min\left\{\sum_{b\in I}z_b:
 0\le z_b\le1,\quad
 \sum_{b\in I}v_p(b)z_b\ge\sum_{a\in A}v_p(a)\ \forall p\right\}.
 \tag{9.2}
\]
Its dual has the form
\[
 L(A,I)=\max_{0\le w_p\le1}
 \left[\sum_{a\in A}X_w(a)-
       \sum_{b\in I}\pos{X_w(b)-1}\right].
 \tag{9.3}
\]
The restriction \(w_p\le1\) follows by clipping a larger coordinate:
the full interval has at least the required total \(p\)-valuation,
so such clipping cannot decrease the dual objective. The bounds
\(z_b\le1\) in (9.2) are essential.

\begin{proposition}[Exact asymptotic criterion]
\label{prop:criterion}
The unrestricted assertion \(g(n)\le(2+o(1))n\) is equivalent to
\[
 \forall\varepsilon>0\ \exists C_\varepsilon<\infty\
 \forall M,x,w:\quad
 \Delta_{2+\varepsilon}(M,x,w)\le C_\varepsilon.
 \tag{9.4}
\]
The constant may depend only on \(\varepsilon\).
\end{proposition}
\begin{proof}
If the asymptotic assertion holds, for each \(\varepsilon>0\) all
cardinalities satisfy
\(g(n)\le(2+\varepsilon)n+C_\varepsilon\), after absorbing finitely
many small \(n\). This also bounds every dual objective.
Maximizing over subsets \(A\subseteq[2,M]\) uses
\[
 \max_A\left(\sum_{a\in A}X_w(a)-(2+\varepsilon)|A|\right)
 =\sum_{a\le M}\pos{X_w(a)-2-\varepsilon}.
 \]
If \(\max A<M\), apply the cover bound inside a contained interval of
length \(\max A\). This proves (9.4).

Conversely, (9.3)--(9.4) give
\(L(A,I)\le(2+\varepsilon)n+C_\varepsilon\).
The all-scale raw fractional transfer of \cite{Wang}, reproduced in
the preceding research note, converts a uniform fractional bound of
this form to the same leading integer coefficient. Explicitly, set
\(\ell(0)=0\) and \(\ell(m)=\sup_{|A|=m}L(A,I)\) for \(m\ge1\). Then
\[
 K_n=\left\lceil\max_{0\le m\le n}
              \{\ell(m)+2(n-m)\}\right\rceil,
 \qquad
 g(n)\le K_n+\pi(2K_n)-1.
 \]
Thus \(g(n)\le(2+\varepsilon)n+C_\varepsilon+o(n)\).
Let \(\varepsilon\) tend to zero.
\end{proof}

The results in this note do \emph{not} establish (9.4) uniformly over
all prime supports. In particular, the following cannot be inferred:

\begin{itemize}
\item A counterexample has been excluded merely because no numerical
search found one.
\item An \(o(M)\) bound for \(\Delta_2\) would suffice. In the difficult
range \(M\asymp n^2\), an error such as \(\sqrt M\) is of order \(n\).
\item Adding one arbitrary prime to a support preserving hinge
dominance permits unlimited iteration. Level dominance is needed
for the next addition and has not been proved in that enlarged class.
\item The rough-source exact theorem proves a weighted hinge theorem
at the same exponent. Weighted sources may contain unweighted
cofactors, so the source classes differ.
\end{itemize}

A remaining counterexample to the exact conjecture must evade the
finite criterion of Section~\ref{sec:exact}. An asymptotic obstruction
must also evade Theorem~\ref{thm:asymclass}, the fixed-support theorem,
and the ordered support and anchor conditions.
These exclusions substantially describe the residual task, but they
do not quantify a percentage of the original problem that is solved.

The particular largest-prime-quotient matching cannot simply be
asserted without a hypothesis. For example, let \(M=440\),
\(I=\{4401,\ldots,4840\}\), and
\[
 A=\{2k: k\le220\text{ is composite},\ k\ne169\}.
 \]
There are 171 sources, while the union of the neighborhoods for
\(a/P^+(a)\) consists of only 170 positions: precisely the interval
multiples of at least one of \(4,6,10,14,22\).
This is a Hall obstruction to that fixed choice of quotient, not a
counterexample to the product-cover problem or to \(H_2\).
An unrestricted proof must accommodate such failures, change the
quotient choices, or use a different argument.

\section{Verification, provenance, and references}
\label{sec:refs}

The fixed-point, ordered-prefix, source-counting, and prime-repair
arguments were independently re-derived within the research session.
Audits specifically checked repeated prime powers, primes outside a
weighted support, nonintegral \(M/d\), arbitrarily large shifts,
empty carrier families, and the requirement that interval positions
remain distinct. The bundle includes exact rational checks and a
finite source-certificate checker. Those checks supplement the
proofs; they do not establish an unrestricted theorem.

The raw LP transfer and the long-interval gate used in Sections
\ref{sec:asymclass} and \ref{sec:remaining} are external inputs.
The exact theorem in Section~\ref{sec:exact} does not depend on those
inputs, on an LP solver, or on the numerical searches.
The standard prime-counting and Mertens estimates are the only
analytic inputs to the exponent in Theorem~\ref{thm:rough}.

\begin{thebibliography}{9}
\bibitem{Problem}
T.~F.~Bloom, \emph{Erd\H{o}s Problem \#708},
\url{https://www.erdosproblems.com/708}.

\bibitem{Chen}
Chen, \emph{Explicit upper bounds for the Erd\H{o}s--Sur\'anyi
function \(g(n)\)}, public working manuscript and verification
materials, including the signed atom-hinge method and long-interval
theorem,
\url{https://github.com/chy4pro/erdos-708-explicit-upper-bounds}.

\bibitem{Wang}
Ziming Wang, \emph{Erd\H{o}s 708 interval divisor covers},
public working manuscript, especially the raw fractional transfer
and the quadratic coefficient-two gate,
\url{https://github.com/science-prince/erdos-708-interval-divisor-covers}.

\bibitem{Bergstrom}
C.~Bergstr\"om, \emph{Scaled-Interval Hinge Telescoping,
Difference-of-Convex Obstructions, and Smooth Proper-Divisor
Matching for Erd\H{o}s Problem 708}, 2026,
\url{https://zenodo.org/records/24366998}.

\bibitem{TaoMertens}
T.~Tao, \emph{Mertens' theorems}, 11 December 2013, especially
the logarithmic product formula (3),
\url{https://terrytao.wordpress.com/2013/12/11/mertens-theorems/}.
\end{thebibliography}
\end{document}
